\section{Density Optimization}\label{sec:density-optimization}

\subsection{Optimize Algorithm}\label{subsec:optimize-algorithm}

\texttt{HyFluid::optimize} operates on one loaded frame set. The frame-set name selects a host record \(H_f\) and the device record \(D_f\) at the same index.

\begin{center}
    \begingroup
    \footnotesize
    \renewcommand{\arraystretch}{1.15}
    \setlength{\tabcolsep}{4pt}
    \begin{tabularx}{\columnwidth}{@{}r>{\raggedright\arraybackslash}X@{}}
        \toprule
        \multicolumn{2}{@{}l@{}}{\textbf{Algorithm 1:} \texttt{HyFluid::optimize}} \\
        \midrule
        1 & Resolve \((H_f,D_f)\) from \texttt{request.frame\_set}. \\
        2 & Repeat for \texttt{request.iterations}. \\
        3 & Sample a training batch from \(D_f\) into ray and sample buffers. \\
        4 & Evaluate the network on sampled coordinates. \\
        5 & Compute image loss and compact valid samples. \\
        6 & Pad the compacted batch to the network granularity. \\
        7 & Re-evaluate the network on compacted samples. \\
        8 & Backpropagate density gradients. \\
        9 & Step the optimizer. \\
        10 & Read counters, update adaptive sample count, increment step. \\
        11 & After the loop, read loss and occupancy counters; return stats. \\
        \bottomrule
    \end{tabularx}
    \endgroup
\end{center}

\begin{center}
    \begingroup
    \scriptsize
    \renewcommand{\arraystretch}{1.12}
    \setlength{\tabcolsep}{3pt}
    \begin{tabularx}{\columnwidth}{@{}p{0.31\columnwidth}p{0.31\columnwidth}>{\raggedright\arraybackslash}X@{}}
        \toprule
        \textbf{State} & \textbf{Used by} & \textbf{Initialized as} \\
        \midrule
        \(H_f\) & sampling dimensions, stats & host frame-set metadata. \\
        \(D_f\) & sampling, loss lookup & uploaded pixels, cameras, intrinsics, frame grid. \\
        \begin{tabular}[t]{@{}l@{}}
            \texttt{device.field\_to\_}\\\texttt{world\_linear}
        \end{tabular}
        & sampling step metric & uploaded field-domain linear map. \\
        \texttt{device.occupancy} & sampling skip mask & initially full occupancy grid. \\
        \texttt{device.params} & network forward/backward & initialized or loaded trainable weights. \\
        \texttt{host.rays\_per\_batch} & sampler batch size & initial training batch shape. \\
        \begin{tabular}[t]{@{}l@{}}
            \texttt{host.inference\_}\\\texttt{sample\_count}
        \end{tabular}
        & network batch size & initialized to max batch, then adapted from sampled count. \\
        \bottomrule
    \end{tabularx}
    \par\smallskip
    \footnotesize\textbf{Table 4:} Initialization anchors for the first optimization algorithm.
    \endgroup
\end{center}
